\section{财富的神经表征}\label{chap3_4:results_wealth_tuning}

前额叶神经元对财富的编码呈现两种不同的调谐模式：单调型和数值型。单调型神经元随财富连续变化，偏好极富或极穷的状态；数值型神经元则对特定财富区间表现出选择性，在偏好区间外放电率下降。如图~\ref{fig4_1:example}A、B所示的ACC神经元放电率随代币数增加单调上升，在最高财富水平达到峰值；图~\ref{fig4_1:example}C、D所示的OFC神经元放电率在中等财富水平达到峰值，在更低或更高财富时均下降。

为量化单神经元对财富状态的编码，我们取选项呈现前$300$至选项呈现后$300$~毫秒的平均放电率对财富水平（$\tilde{W}_5$，五水平类别变量，细节见~\ref{chap2:methods_neural_wealth}~节）进行线性回归。神经元的偏好财富水平分布呈U型，偏好极低和极高财富状态的神经元比例略高（图~\ref{fig4_2:proportion}）。

群体水平的调谐曲线进一步刻画了财富编码的整体形态（图~\ref{fig4_3:tuning}）。将神经元按偏好财富水平分组取平均后，各组曲线呈现相似的峰值结构。将所有神经元的财富调谐曲线以偏好财富水平为中心重新对齐后，±1与±2位置的归一化放电率之间存在显著差异，表明调谐曲线在峰值两侧呈现真实的梯度衰减，而非随机噪声。打乱财富组标签后该差异消失，确认了群体调谐结构的真实性。

\begin{figure}[H]
    \centering
    \includegraphics[width=0.85\textwidth]{Figure4_1_Example}
    \bicaption{\enspace 编码财富的示例神经元}{\enspace Representative neurons encoding wealth}
    \label{fig4_1:example}
\end{figure}
\fignotetext{\textbf{(A) 示例ACC神经元（单调型）。}选项呈现前后$[-300, 300]$~毫秒窗口内的放电率（均值$\pm$标准误）。试次按财富水平分为五组（[1]、[2–3]、[4–7]、[8–15]、[16–24]）以确保各组试次数大致均衡。颜色由深至浅表示财富水平递增。\\ \textbf{(B) 图A神经元的调谐曲线。}纵轴为分析窗口内各组平均放电率，横轴为财富水平分组；该神经元在最高财富组（[16–24]）达到峰值。\\ \textbf{(C) 示例OFC神经元（数值型）。}布局同(A)。\\ \textbf{(D) 图C神经元的调谐曲线。}布局同(B)；该神经元在[4–7]组达到峰值，更低或更高财富时放电率均下降。}{\textbf{(A) A representative ACC neuron (monotonic tuning).} Firing rate (mean $\pm$ SEM) in the $[-300, 300]$~ms window aligned to offer onset. Trials are grouped into five wealth levels ([1], [2–3], [4–7], [8–15], [16–24]) to ensure roughly equal trial counts across groups; darker to lighter colors indicate increasing wealth. \\ \textbf{(B) Tuning curve of the neuron in (A).} The y-axis shows the mean firing rate within the analysis window for each wealth group; this neuron peaks at the highest wealth group ([16–24]). \\ \textbf{(C) A representative OFC neuron (numeric tuning).} Layout as in (A). \\ \textbf{(D) Tuning curve of the neuron in (C).} Layout as in (B); this neuron peaks at the [4–7] group and shows reduced firing at both lower and higher wealth levels.}

\begin{figure}[H]
    \centering
    \includegraphics[width=0.85\textwidth]{Figure4_2_Proportion}
    \bicaption{\enspace 神经元偏好的财富水平分布}{\enspace Distribution of preferred wealth level}
    \label{fig4_2:proportion}
\end{figure}
\fignotetext{各脑区中偏好每个财富水平的财富编码神经元比例。第一行为三只猕猴的聚合数据，第二至四行分别为猕猴W、U、T的数据；每种颜色对应一个脑区（ACC：红色、DLPFC：绿色、OFC：蓝色）。财富编码神经元定义为在$[-300, 300]$~毫秒窗口内放电率对$\tilde{W}_5$回归的F检验显著（$p < 0.05$）的神经元；偏好财富水平取五组比较中$p$值最小的水平（详见~\ref{chap2:methods_neural_wealth_tuning}）。彩色条为实际数据，灰色条为随机水平。}{Proportion of wealth-encoding neurons preferring each wealth level, shown for each brain area. Row 1 shows data pooled across all three monkeys; rows 2–4 show data for Monkey W, U, and T, respectively. Different colors correspond to different brain areas (ACC: red, DLPFC: green, OFC: blue). A neuron is classified as wealth-encoding if the F-test for $\tilde{W}_5$ in the $[-300, 300]$~ms regression is significant ($p < 0.05$); the preferred wealth level is defined as the level yielding the smallest $p$-value in one-versus-all pairwise comparisons (see Section~\ref{chap2:methods_neural_wealth_tuning}). Colored bars show the actual data; gray bars show the shuffle control. The distribution of preferred wealth levels is U-shaped in all three areas, and this pattern is consistent across individual monkeys.}

\begin{figure}[H]
    \centering
    \includegraphics[width=0.9\textwidth]{Figure4_3_TuningCurve}
    \bicaption{\enspace 财富水平调谐曲线}{\enspace Tuning curves for wealth level}
    \label{fig4_3:tuning}
\end{figure}
\fignotetext{四个子图分别展示聚合数据（左上）、猕猴W（右上）、猕猴U（左下）、猕猴T（右下）。每个子图内左列为按偏好财富水平分组的归一化平均调谐曲线，红、绿、蓝分别对应ACC、DLPFC、OFC，颜色由深至浅表示偏好财富水平从低（[1]）到高（[16–24]）；右列为将各组曲线以其偏好财富水平为中心重新对齐后的结果。聚合数据中，重新对齐后±1与±2位置的归一化放电率存在显著差异（ACC：$p = 3.31\times10^{-41}$；DLPFC：$p = 2.50\times10^{-36}$；OFC：$p = 2.32\times10^{-36}$；Wilcoxon秩和检验）。打乱财富标签后重复相同分析，±1与±2位置差异不再显著（ACC：$p = 0.313$；DLPFC：$p = 0.251$；OFC：$p = 0.164$），个体结果与聚合结果一致。}{Four panels show data for the pooled sample (top left), Monkey W (top right), Monkey U (bottom left), and Monkey T (bottom right), respectively. Within each panel, the left column shows normalized mean tuning curves grouped by preferred wealth level; red, green, and blue correspond to ACC, DLPFC, and OFC, respectively, with darker to lighter shades indicating preferred wealth levels from low ([1]) to high ([16–24]). The right column shows tuning curves re-centered on each neuron's preferred wealth level, revealing a gradient of decreasing firing rate with distance from the peak. In the pooled data, normalized firing rates at $\pm1$ and $\pm2$ positions relative to the peak differed significantly (ACC: $p = 3.31\times10^{-41}$; DLPFC: $p = 2.50\times10^{-36}$; OFC: $p = 2.32\times10^{-36}$; Wilcoxon rank-sum test). After permuting wealth group labels, this difference was no longer significant (ACC: $p = 0.313$; DLPFC: $p = 0.251$; OFC: $p = 0.164$), and individual monkey results were consistent with the pooled analysis.}

\begin{figure}[H]
    \centering
    \includegraphics[width=0.9\textwidth]{Figure4_3_2_Skew}
    \bicaption{\enspace 财富调谐曲线的偏态参数分布}{\enspace Distribution of skewness parameters of wealth tuning curves}
    \label{fig4_3_2:skew}
\end{figure}
\fignotetext{图中展示在线性财富坐标（珊瑚色）和对数财富坐标（蓝色）下，对各神经元调谐曲线拟合Skew-Normal分布所得偏态参数的分布直方图。布局为4行$\times$3列，行依次对应聚合数据、猕猴W、猕猴U、猕猴T，列依次对应ACC、DLPFC、OFC。每个子图左上角标注神经元数量（$n$），灰色虚线为偏态参数$=0$参考线，珊瑚色和蓝色数显分别为线性和对数坐标下的中位数，若显著偏离0则用实线标注，否则为虚线。聚合数据中，DLPFC和OFC在对数坐标下偏态参数与0无显著差异（DLPFC：$p_\text{log}=8.1\times10^{-1}$；OFC：$p_\text{log}=8.9\times10^{-1}$），在线性坐标下显著大于0（DLPFC：$p_\text{lin}=2.6\times10^{-13}$；OFC：$p_\text{lin}=5.1\times10^{-11}$）；ACC在两种坐标下均显著大于0（$p_\text{lin}=2.6\times10^{-14}$，$p_\text{log}=3.4\times10^{-7}$）。}{Histograms show the distribution of the skewness parameter from skew-normal fits to each neuron's tuning curve, under linear (coral) and logarithmic (blue) wealth coordinates. The layout is 4 rows $\times$ 3 columns, with rows corresponding to pooled data, Monkey W, Monkey U, and Monkey T, and columns corresponding to ACC, DLPFC, and OFC. Within each panel, the top-left annotation reports the number of neurons ($n$); the gray dashed line marks skewness $= 0$; the coral and blue lines indicate the median skewness under linear and logarithmic coordinates, respectively, with solid ones indicating significant deviation from 0 and dashed lines otherwise. In the pooled data, skewness did not differ significantly from 0 under logarithmic coordinates in DLPFC and OFC (DLPFC: $p_\text{log} = 8.1\times10^{-1}$; OFC: $p_\text{log} = 8.9\times10^{-1}$), but was significantly positive under linear coordinates (DLPFC: $p_\text{lin} = 2.6\times10^{-13}$; OFC: $p_\text{lin} = 5.1\times10^{-11}$). ACC showed significantly positive skewness under both coordinate systems ($p_\text{lin} = 2.6\times10^{-14}$; $p_\text{log} = 3.4\times10^{-7}$).}

调谐曲线的对称性可以进一步揭示财富状态编码的内在尺度。Nieder等人（2003）发现，猕猴前额叶皮层对数量的神经编码符合Weber-Fechner定律：在对数坐标下调谐曲线呈对称分布，而在线性坐标下则表现为右偏。本实验中的财富状态用代币数量表示，与Nieder等人的实验设置十分相似。用偏态高斯分布拟合财富调谐曲线发现，DLPFC和OFC在对数坐标下偏态参数与0无显著差异，在线性坐标下偏态参数则显著大于0（图~\ref{fig4_3_2:skew}），说明这两个脑区对财富的神经表征也符合Weber-Fechner定律。ACC的结果则有所不同，在两种坐标下偏态参数均显著大于0，这说明ACC可能以超对数方式压缩压缩编码财富水平。

以上结果表明，前额叶神经元通过异质但有序的调谐，在群体水平上实现了对财富状态的压缩表征。